% !TeX root = ../../Book2.tex
% 纲要标题：### A.1 整数与一元多项式矩阵的 Smith 标准形算法
% 纲要位置：计划纲要.md，第 1744 行。

% 纲要要点（供目录核对）：
%
% * 整数与一元多项式矩阵的 Smith 标准形及精确算法；保持矩阵等式与可逆性，区分精确商清除和余数降主元，以局部下降及有限主元阶段证明终止；正式证明整数证书的三个检查条件及原坐标的解、核、余核
% #### A.1.1 行列变换的记录
% #### A.1.2 主元约化与整除链
% #### A.1.3 标准形证书与线性方程
%

\section{整数与一元多项式矩阵的 Smith 标准形算法}
\label{b2:sec:A01}

\chapref{b2:ch:04}已经证明主理想整环上 Smith 标准形的存在性和唯一性。实际求解时，还要回答两个问题：允许哪些运算，怎样确认最终输出确实对应原来的矩阵？本节在整数环或有效给定的域上的一元多项式环中使用带余除法，记录每一步换基，使标准形同时给出原坐标中的解、核与余核。

\subsection{行列变换的记录}
\label{b2:subsec:A01:01}

设 \(R=\mathbb Z\) 或 \(R=k[t]\)，其中域 \(k\) 的加减乘除和相等判定均可精确执行。输入 \(A\in M_{m\times n}(R)\)，从
\[
B=A,\qquad P=I_m,\qquad Q=I_n
\]
出发，每一步都保持 \(B=PAQ\)、\(P\in\operatorname{GL}_m(R)\) 和 \(Q\in\operatorname{GL}_n(R)\)。作行变换 \(B\leftarrow EB\) 时同步令 \(P\leftarrow EP\)，因为 \(EB=(EP)AQ\)；作列变换 \(B\leftarrow BF\) 时同步令 \(Q\leftarrow QF\)，因为 \(BF=PA(QF)\)。这里 \(E,F\) 分别是相应尺寸的可逆初等矩阵，所以更新后的 \(P,Q\) 仍可逆。这些等式和可逆性是算法始终保持的不变量；行变换的矩阵累积在 \(P\) 的左侧，列变换的矩阵累积在 \(Q\) 的右侧。

\ProseParagraphBreak
这些规则看似只是在保存计算过程，却决定了最后能否把解送回原坐标。若 \(x=Qy\)，原方程 \(Ax=b\) 恰好变成 \(By=Pb\)。因此右端要随行变换改变，而未知量的坐标要用 \(Q\) 变回去，不能把两个换基矩阵混用。

\begin{example}\label{b2:exm:A-integer-certificate}
取整数矩阵
\[
A=\begin{pmatrix}4&6&2\\2&6&4\end{pmatrix}.
\]
先交换两行，再从第二行减去第一行的两倍，最后把第二行乘 \(-1\)，得到
\[
\begin{pmatrix}2&6&4\\0&6&6\end{pmatrix}.
\]
接着依次作 \(C_2\leftarrow C_2-3C_1\)、\(C_3\leftarrow C_3-2C_1\)、\(C_3\leftarrow C_3-C_2\)，其中最后一步的 \(C_2\) 指已经更新的列。结果及累积矩阵为
\begin{equation}\label{b2:eq:A-integer-certificate}
P=\begin{pmatrix}0&1\\-1&2\end{pmatrix},\qquad
Q=\begin{pmatrix}1&-3&1\\0&1&-1\\0&0&1\end{pmatrix},\qquad
PAQ=\begin{pmatrix}2&0&0\\0&6&0\end{pmatrix}.
\end{equation}
两换基矩阵的行列式均为一，且 \(2\mid6\)，所以这个等式是 Smith 标准形的完整证书。原矩阵的一阶子式最大公因子为二，三个二阶子式分别为 \(12,12,12\)，也与对角元的乘积相符。
\end{example}

\subsection{主元约化与整除链}
\label{b2:subsec:A01:02}

整数带余除法使用非负余数，多项式带余除法使用次数严格较小的余式。对非零主元 \(a\)，若同一列中有 \(b\)，写 \(b=qa+r\)，减去相应的 \(q\) 倍主元行。余数非零时交换这两行，新的主元是 \(r\)；余数为零时，\(q=b/a\) 是环中的精确商，行减法直接把该项清成零。处理同行元素则使用列变换。\cref{b2:alg:A-smith}中的 \(j\) 表示当前尚未处理的首个行列位置，每次只在编号至少为 \(j\) 的行、列中操作，保留此前已经完成的对角块。

\begin{algorithm}[htbp]
\caption{Smith 约化}
\label{b2:alg:A-smith}
\begin{algorithmsteps}{\(A\in M_{m\times n}(R)\)，其中 \(R=\mathbb Z\) 或 \(k[t]\)，\(k\) 为有效给定的域}{幺模矩阵 \(P,Q\) 及满足 \(PAQ=D\) 的 Smith 标准形 \(D\)}
\State 初始化 \(B=A\)、\(P=I_m\)、\(Q=I_n\)、\(j=1\)。行变换同步令 \(P\leftarrow EP\)，列变换同步令 \(Q\leftarrow QF\)，保持 \(B=PAQ\) 及 \(P,Q\) 可逆。
\While{\(j\le\min(m,n)\)}
\State 若右下剩余子块为零，结束；否则把其中一个非零项移到 \((j,j)\)。
\State 用带余除法检查第 \(j\) 列下方和第 \(j\) 行右方。发现非零余数就把它换成主元，并重新检查；全被主元整除时，先用精确商作行减法清除该列下方，再作列减法清除该行右方。
\State 若主元 \(a=B_{jj}\) 不整除右下子块某项 \(b=B_{ik}\)，把第 \(i\) 行加到第 \(j\) 行；对 \(b=qa+r\) 作 \(C_k\leftarrow C_k-qC_j\)，交换第 \(j,k\) 列，使非零余数 \(r\) 成为新主元，然后返回上一步。否则保留当前主元，令 \(j\leftarrow j+1\)。
\EndWhile
\State 把非零对角元归一化为正整数或首一多项式，仍同步更新换基矩阵；输出 \(P,Q,B\)。
\end{algorithmsteps}
\end{algorithm}

分块检查也确实会缩小主元。清除第 \(j\) 行、列后，若右下子块的 \((i,k)\) 项 \(b\) 不被 \(a=B_{jj}\) 整除，只看第 \(j,i\) 行和第 \(j,k\) 列，相关部分的变换为
\[
\begin{pmatrix}a&0\\0&b\end{pmatrix}
\xrightarrow{R_j\leftarrow R_j+R_i}
\begin{pmatrix}a&b\\0&b\end{pmatrix}
\xrightarrow{C_k\leftarrow C_k-qC_j}
\begin{pmatrix}a&r\\0&b\end{pmatrix},
\qquad b=qa+r,\quad r\ne0.
\]
交换这两列后，\(r\) 成为新主元。上式只画出有关的两行、两列；把第 \(i\) 行加到第 \(j\) 行，也会把第 \(i\) 行的其他条目带入第 \(j\) 行，所以必须返回行列检查，重新清除。这里暂时破坏的是当前行、列的零项，目的是取得更小主元；编号小于 \(j\) 的已完成块不受影响，因为剩余行在这些列上为零，剩余列在这些行上也为零。

\ProseParagraphBreak
固定 \(j\) 后，每次遇到非零余数都会严格缩小主元：整数情形为 \(0<r<|a|\)，多项式情形为 \(\deg r<\deg a\)。右下块的整除障碍经上述变换也造成同样的下降，因而主元只能更换有限次。在两次更换之间，只需检查有限个条目；若余数全部为零，先清列再清行的精确商操作也只有有限次，且清行不会重新产生已清掉的列下方条目。最终主元整除整个剩余子块，\(j\) 才增加。由于 \(j\le\min(m,n)\)，整个算法也在有限步后结束。

\ProseParagraphBreak
保留当前主元后，后续只在右下子块内作行列变换，新条目仍是旧条目的 \(R\) 线性组合，所以仍被这个主元整除。这保证后续得到的各个非零对角元接成整除链。最后的归一化也只乘单位：负整数主元所在的行乘 \(-1\)，多项式主元若首项系数为 \(c\ne0\)，所在的行乘 \(c^{-1}\)。因此正值或首一归一化不改变幺模等价，仍可同步保留完整的 \(P,Q\)。

\ProseParagraphBreak
这给出在 \(\mathbb Z\) 与 \(k[t]\) 上实现\cref{b2:thm:smith-normal-form}的算法，终止依据是可计算的绝对值或次数下降。一般主理想整环上的存在证明使用 Bézout 变换与主元理想的升链条件；要执行相应计算，还需具备\secref{b2:sec:A02}列出的有效运算。只做到对角化仍不够：例如 \(\operatorname{diag}(6,10)\) 不满足 Smith 条件，因为 \(6\nmid10\)；其一阶子式最大公因子为二，行列式为六十，\cref{b2:thm:smith-uniqueness}给出的归一化对角元应为 \(2,30\)。

\subsection{标准形证书与线性方程}
\label{b2:subsec:A01:03}

\begin{proposition}\label{b2:prop:A-smith-solving}
设 \(R\) 是主理想整环，\(A\in M_{m\times n}(R)\)，并已知
\[
PAQ=D,\qquad P\in\operatorname{GL}_m(R),\quad Q\in\operatorname{GL}_n(R),
\]
其中 \(D\) 是非零对角元为 \(d_1\mid\cdots\mid d_r\) 的 Smith 标准形。给定 \(b\in R^m\)，记 \(c=Pb\)。方程 \(Ax=b\) 有解当且仅当
\[
d_i\mid c_i\quad(1\le i\le r),\qquad c_i=0\quad(r<i\le m).
\]
有解时，全部解是 \(x=Qy\)，其中 \(y_i=c_i/d_i\) 对 \(i\le r\) 成立，其余 \(n-r\) 个坐标任取。特别地，若 \(e_i\) 是 \(R^n\) 的第 \(i\) 个标准基向量，则 \(\ker A\) 的一组基为 \(Qe_{r+1},\ldots,Qe_n\)。此外，有余核分解
\[
R^m/A(R^n)\cong\bigoplus_{i=1}^r R/(d_i)\ \oplus R^{m-r}.
\]
若 \(f_i\) 是 \(R^m\) 的第 \(i\) 个标准基向量，则第 \(i\) 个分量的坐标生成元在原余核中由 \(P^{-1}f_i+A(R^n)\) 代表：\(i\le r\) 时对应 \(R/(d_i)\)，\(i>r\) 时对应自由分量。
\end{proposition}
\begin{proof}
代入 \(x=Qy\) 并左乘 \(P\)，把原方程等价变为 \(Dy=c\)。前 \(r\) 行分别是 \(d_i y_i=c_i\)，其可解性就是整除条件；由于 \(d_i\ne0\) 且 \(R\) 为整环，可解时该坐标唯一。剩余行左端为零，正好给出所述零坐标条件；剩余列不受限制。\(Q\) 是同构，把这些零列对应的自由坐标轴送成原核的一组基。

由\cref{b2:prop:smith-cokernel}，余核同构把 \(v+A(R^n)\) 送到 \(Pv+D(R^n)\)；其逆把对角坐标 \(f_i\) 的类送到 \(P^{-1}f_i+A(R^n)\)。前 \(r\) 个坐标模掉 \((d_i)\)，余下零行的坐标则没有关系约束，所以得到所述循环分量与自由分量。若 \(d_i\) 是单位，对应循环分量为零。
\end{proof}

在\cref{b2:exm:A-integer-certificate}中取 \(b=(10,8)^{\mathsf T}\)，则 \(Pb=(8,6)^{\mathsf T}\)，所以
\[
x=\begin{pmatrix}1\\1\\0\end{pmatrix}
  +z\begin{pmatrix}1\\-1\\1\end{pmatrix},\qquad z\in\mathbb Z.
\]
直接乘原矩阵分别得到 \(b\) 和零，便同时核验了一个特解与齐次解。若把右端换成 \((8,8)^{\mathsf T}\)，则变换后的第二个坐标是八，不能被六整除，故没有整数解，即使在 \(\mathbb Q\) 上仍有解。

\ProseParagraphBreak
零列给出核的自由坐标，零行给出余核的自由坐标，两者在矩形矩阵中不能混同。若系数环有零因子，非零对角元的零化子也可能贡献核，不能只读取零列；\cref{b2:prop:A-diagonal-annihilators}将给出相应的一般公式。若只要抽象余核的同构类型，Smith 对角数据已经足够；恢复余核生成元的原坐标使用 \(P^{-1}\)，恢复未知量或核的原坐标则使用 \(Q\)。

\begin{proposition}\label{b2:prop:A-smith-certificate-validation}
设 \(m,n\ge1\) 为整数，\(A,D\in M_{m\times n}(\mathbb Z)\)，\(P\in M_m(\mathbb Z)\)、\(Q\in M_n(\mathbb Z)\)。若
\[
PAQ=D,\qquad \det P,\det Q\in\{1,-1\},
\]
且 \(D\) 是矩形对角矩阵，其非零对角元占据前 \(r\) 个位置，其中整数 \(r\) 与这些对角元满足
\[
0\le r\le\min(m,n),\qquad
d_i>0\ (1\le i\le r),\qquad d_1\mid\cdots\mid d_r,
\]
其余条目均为零，则 \((P,Q,D)\) 给出 \(A\) 的正归一化整数 Smith 标准形证书。\(D\) 由 \(A\) 唯一确定，而换基矩阵 \(P,Q\) 不要求唯一。验证这样的证书必须同时检查整数矩阵等式、整数上的可逆性以及归一化整除链。
\end{proposition}
\begin{proof}
整数矩阵的行列式为 \(\pm1\) 时，伴随矩阵公式给出的逆矩阵仍有整数条目，所以 \(P\in\operatorname{GL}_m(\mathbb Z)\)、\(Q\in\operatorname{GL}_n(\mathbb Z)\)。因此所给矩阵等式证明 \(A,D\) 在整数环上幺模等价。对角形、非零元排列和整除链满足\cref{b2:def:smith-normal-form}的全部条件，故 \(D\) 确为 Smith 标准形。\cref{b2:thm:smith-uniqueness}确定各非零对角元的相伴类，正值归一化再确定唯一代表。反过来，整数幺模矩阵的行列式必为 \(\pm1\)，而正归一化的 Smith 形必满足所写排列及整除条件；这些检查与计算标准形时采取了哪条变换路径无关。

例如取 \(A=D=\operatorname{diag}(2,3)\)、\(P=Q=I_2\)，矩阵等式与整数可逆性均成立，但 \(2\nmid3\)，所以这不是 Smith 证书。其条目最大公因子为一，行列式为六，由\cref{b2:thm:smith-uniqueness}，正归一化 Smith 形应为 \(\operatorname{diag}(1,6)\)。具体取
\[
P_0=\begin{pmatrix}1&1\\-3&-2\end{pmatrix},\qquad
Q_0=\begin{pmatrix}-1&-3\\1&2\end{pmatrix},\qquad
P_0\operatorname{diag}(2,3)Q_0=\operatorname{diag}(1,6).
\]
\(Q_0\) 的第一列使原对角矩阵的像为 \((-2,3)^{\mathsf T}\)，再由 \(P_0\) 的第一行读取 Bézout 等式 \(-2+3=1\)，取得单位主元；\(Q_0\) 的第二列使这一行的另一项为零。两换基矩阵的行列式都为一，所以这是合法的整数证书。另一方面，有理矩阵等式
\[
\operatorname{diag}(1/2,1/3)\operatorname{diag}(2,3)=I_2
\]
也不能给出整数 Smith 证书：左侧缩放矩阵根本不是整数矩阵，更不属于 \(\operatorname{GL}_2(\mathbb Z)\)。事实上，\(\operatorname{diag}(2,3)\) 的整数余核为 \(\mathbb Z/2\mathbb Z\oplus\mathbb Z/3\mathbb Z\ne0\)，而 \(I_2\) 的整数余核为零，故二者不可能在整数环上幺模等价。
\end{proof}
